%=====================================================================
\chapter{What Machine Learning Is}\label{ch:introduction}
%=====================================================================
\locator{1}{Part I --- Foundations}

\begin{outcomes}
By the end of this chapter you will be able to:
\begin{tight}
  \item \textbf{Define} machine learning precisely, and name the task,
        performance measure and experience of a given learning problem.
  \item \textbf{Distinguish} learning a program from writing one, and machine
        learning from the wider field of artificial intelligence.
  \item \textbf{Classify} a problem as supervised (regression or
        classification), unsupervised, semi-supervised or reinforcement learning.
  \item \textbf{Describe} the four design decisions in building a learning system.
  \item \textbf{Explain} why a learner must assume something the data does not
        say in order to generalise at all.
\end{tight}
\end{outcomes}

\begin{prereq}
The course \emph{Artificial Intelligence}, in which you met agents, search and
reasoning under uncertainty. Nothing in this handout is assumed yet.
\end{prereq}

\begin{keyterms}
machine learning $\bullet$ task, performance measure, experience $\bullet$
training data $\bullet$ feature $\bullet$ label $\bullet$ model $\bullet$
supervised learning $\bullet$ classification $\bullet$ regression $\bullet$
unsupervised learning $\bullet$ semi-supervised learning $\bullet$
reinforcement learning $\bullet$ target function $\bullet$ hypothesis
$\bullet$ hypothesis space $\bullet$ generalisation $\bullet$ inductive bias
\end{keyterms}

\section{Why Learn a Program Instead of Writing One?}

Suppose you are asked to write a program that flags phishing e-mails. You could
sit down and write rules: \emph{if the sender's domain does not match the display
name, and the message contains a link, and it mentions a password, flag it.}
The first version catches some phishing. Then the attackers read your rules ---
or simply vary their wording --- and the rules stop working. You add more. After
a year you have four hundred rules that contradict one another, nobody
understands, and that still miss the attack that arrived this morning.

Now suppose instead that you collect fifty thousand e-mails that people have
already marked as phishing or legitimate, and hand them to a program whose job
is to \emph{find} the rules. When the attackers change tactics, you do not
rewrite anything: you collect new examples and let the program find new rules.
That is the whole idea of machine learning, and \dsref{program-vs-learn} shows
it in one picture.

\dsfig{%
\begin{tikzpicture}[font=\scriptsize,
  io/.style={rounded corners=3pt, draw=#1, line width=0.9pt, fill=#1!8,
             text=#1!55!black, font=\scriptsize\bfseries, minimum width=1.9cm,
             minimum height=0.7cm, align=center},
  engine/.style={rounded corners=4pt, fill=#1, text=white, font=\scriptsize\bfseries,
                 minimum width=2.7cm, minimum height=1.25cm, align=center},
  a/.style={-{Stealth[length=2.2mm]}, gray!60, line width=0.9pt}]
% --- traditional programming
\node[font=\small\bfseries, text=primary, anchor=west] at (-1.1,2.35) {Traditional programming};
\node[io=primary] (d1) at (0,1.5)  {Data};
\node[io=accent]  (r1) at (0,0.55) {Rules\\{\tiny written by you}};
\node[engine=primary] (c1) at (3.6,1.02) {Computer\\runs the rules};
\node[io=greenhl] (a1) at (7.1,1.02) {Answers};
\draw[a] (d1) -- (c1);
\draw[a] (r1) -- (c1);
\draw[a] (c1) -- (a1);
% --- machine learning
\node[font=\small\bfseries, text=greenhl!55!black, anchor=west] at (-1.1,-0.55) {Machine learning};
\node[io=primary] (d2) at (0,-1.4)  {Data};
\node[io=greenhl] (a2) at (0,-2.35) {Answers\\{\tiny (labels)}};
\node[engine=greenhl!70!black] (c2) at (3.6,-1.88) {Learning\\algorithm};
\node[io=accent]  (r2) at (7.1,-1.88) {Rules\\{\tiny = the model}};
\draw[a] (d2) -- (c2);
\draw[a] (a2) -- (c2);
\draw[a] (c2) -- (r2);
% --- then the model is used like a program
\node[io=primary] (d3) at (9.9,-0.55) {New data};
\node[io=greenhl] (a3) at (9.9,-3.2) {Predicted\\answers};
\draw[a] (d3) -- (r2);
\draw[a] (r2) -- (a3);
\node[anchor=north west, align=left, text width=4.6cm, font=\scriptsize, text=gray!55!black]
     at (8.3,2.45)
     {\textbf{The swap.} In traditional programming the rules are an \emph{input};
      in machine learning they are the \emph{output}. You supply examples of the
      answers you want, and the algorithm writes the rules.};
\end{tikzpicture}}%
{Traditional programming against machine learning. The learned model is then used
exactly like a program: new data in, answers out.}{program-vs-learn}

\begin{conceptbox}[When Learning Beats Writing]
Machine learning earns its cost in four situations: when the rules are too
numerous to write (recognising handwriting); when nobody can articulate them
(judging whether a machine sounds unhealthy); when they change faster than a
person can rewrite them (phishing, fraud); and when they must be different for
every user (recommendations). If you can write the rule down in an afternoon and
it will not change, write it --- a rule is cheaper, faster and easier to explain
than any model.
\end{conceptbox}

\section{A Precise Definition}

The definition used throughout this course is Tom Mitchell's, from the set text.
It is worth memorising because it forces you to say exactly what you mean before
you start.

\begin{definitionbox}[Machine Learning (Mitchell, 1997)]
A computer program is said to \term{learn} from experience $E$ with respect to
some class of tasks $T$ and performance measure $P$, if its performance at tasks
in $T$, as measured by $P$, improves with experience $E$.
\end{definitionbox}

The three letters are not decoration. A project that cannot name its $P$ cannot
tell whether it is making progress; a project that cannot name its $E$ does not
know what data to collect. Most failed machine learning projects failed here,
before any algorithm was chosen.

\begin{worked}[Naming T, P and E]
State the task, performance measure and experience for three problems.

\smallskip
\textbf{1. Phishing filter.}
\begin{tight}
  \item $T$: classify each incoming e-mail as phishing or legitimate.
  \item $P$: the fraction of phishing e-mails caught, subject to fewer than one
        legitimate e-mail in a thousand being blocked.
  \item $E$: a corpus of e-mails already labelled by users and analysts.
\end{tight}

\textbf{2. Sprint effort estimation.}
\begin{tight}
  \item $T$: predict the hours a user story will take.
  \item $P$: average absolute error, in hours, on stories not yet estimated.
  \item $E$: past stories with their descriptions, story points and actual hours.
\end{tight}

\textbf{3. A robot learning to balance.}
\begin{tight}
  \item $T$: choose motor commands that keep the robot upright.
  \item $P$: average time before it falls.
  \item $E$: the robot's own attempts --- no labelled examples at all, only the
        consequences of what it tried.
\end{tight}

\smallskip
\textbf{Notice the choice of $P$ in problem 1.} ``Accuracy'' would have been the
obvious answer and the wrong one: if 2\% of e-mail is phishing, a filter that lets
everything through is 98\% accurate. Choosing $P$ is a design decision with
consequences, and \chref{evaluation} is devoted to it.
\end{worked}

\section{The Four Paradigms}

What varies most between learning problems is the \emph{experience} --- in
particular, what the data tells you about the right answer. That gives four
paradigms, sketched in \dsref{paradigms}.

\dsfig{%
\begin{tikzpicture}[font=\scriptsize,
  panel/.style={rounded corners=4pt, draw=#1, line width=0.9pt, fill=#1!4,
                minimum width=6.6cm, minimum height=4.0cm},
  ttl/.style={font=\small\bfseries, text=#1!65!black, anchor=north west},
  sub/.style={font=\scriptsize, text=gray!55!black, anchor=north west, align=left,
              text width=6.2cm},
  ptp/.style={circle, fill=secondary, inner sep=1.5pt},
  ptn/.style={rectangle, fill=accent, inner sep=1.7pt},
  unl/.style={circle, fill=gray!45, inner sep=1.4pt}]
% ---------- supervised
\node[panel=primary] at (0,0) {};
\node[ttl=primary] at (-3.2,1.9) {Supervised};
\node[sub] at (-3.2,1.45) {every example comes with its answer};
\foreach \x/\y in {-2.3/-0.2,-1.8/0.4,-1.3/-0.5,-2.0/-1.0,-1.1/0.1,-2.6/0.6}{
  \node[ptp] at (\x,\y) {};}
\foreach \x/\y in {0.4/-0.3,0.9/0.5,1.3/-0.8,0.6/-1.2,1.7/0.1,2.2/-0.4}{
  \node[ptn] at (\x,\y) {};}
\draw[primary, line width=1.1pt, dashed] (-0.55,0.8) -- (-0.15,-1.6);
\node[font=\tiny, text=primary, anchor=west] at (-0.05,-1.55) {learned boundary};
% ---------- unsupervised
\node[panel=purplehl] at (7.0,0) {};
\node[ttl=purplehl] at (3.8,1.9) {Unsupervised};
\node[sub] at (3.8,1.45) {no answers at all --- find the structure};
\foreach \x/\y in {5.0/-0.1,5.4/0.4,5.7/-0.4,5.2/-0.8}{\node[unl] at (\x,\y) {};}
\foreach \x/\y in {8.2/0.5,8.7/0.2,8.4/-0.2,8.9/0.7}{\node[unl] at (\x,\y) {};}
\foreach \x/\y in {7.4/-1.3,7.9/-1.1,7.6/-0.9}{\node[unl] at (\x,\y) {};}
\draw[purplehl, line width=0.9pt] (5.35,-0.22) ellipse (0.7cm and 0.85cm);
\draw[purplehl, line width=0.9pt] (8.55,0.3) ellipse (0.72cm and 0.6cm);
\draw[purplehl, line width=0.9pt] (7.63,-1.1) ellipse (0.55cm and 0.4cm);
% ---------- semi-supervised
\node[panel=tealhl] at (0,-4.4) {};
\node[ttl=tealhl] at (-3.2,-2.5) {Semi-supervised};
\node[sub] at (-3.2,-2.95) {a few answers, many examples without};
\foreach \x/\y in {-2.5/-4.3,-2.1/-3.8,-1.7/-4.7,-2.3/-5.3,-1.4/-4.2,-2.8/-4.8,
                   -1.9/-5.0,-1.2/-3.6}{\node[unl] at (\x,\y) {};}
\foreach \x/\y in {0.6/-4.1,1.0/-4.9,1.5/-3.8,0.8/-5.4,1.9/-4.5,2.3/-5.1,1.3/-4.4,
                   2.1/-3.7}{\node[unl] at (\x,\y) {};}
\node[ptp, minimum size=5pt] at (-2.1,-4.5) {};
\node[ptn, minimum size=5pt] at (1.5,-4.8) {};
\draw[tealhl, line width=1.1pt, dashed] (-0.4,-3.5) -- (-0.3,-5.9);
\node[font=\tiny, text=tealhl!70!black, anchor=north west, align=left] at (-0.2,-5.35)
     {two labels, and the gap\\between the groups};
% ---------- reinforcement
\node[panel=goldhl] at (7.0,-4.4) {};
\node[ttl=goldhl] at (3.8,-2.5) {Reinforcement};
\node[sub] at (3.8,-2.95) {no answers --- only rewards for what was tried};
\node[rounded corners=3pt, fill=goldhl, text=white, font=\scriptsize\bfseries,
      minimum width=1.6cm, minimum height=0.7cm] (ag) at (5.3,-4.8) {Agent};
\node[rounded corners=3pt, fill=goldhl!25, draw=goldhl, text=goldhl!40!black,
      font=\scriptsize\bfseries, minimum width=1.9cm, minimum height=0.7cm] (en)
      at (8.7,-4.8) {Environment};
\draw[-{Stealth[length=2mm]}, goldhl!70!black, line width=0.9pt]
     (ag.north) .. controls +(0,0.75) and +(0,0.75) .. node[above, font=\tiny] {action} (en.north);
\draw[-{Stealth[length=2mm]}, goldhl!70!black, line width=0.9pt]
     (en.south) .. controls +(0,-0.75) and +(0,-0.75) .. node[below, font=\tiny]
     {new state, reward} (ag.south);
\end{tikzpicture}}%
{The four paradigms, distinguished by what the experience tells the learner about the
right answer: all of it, none of it, a little of it, or only whether what it did turned
out well.}{paradigms}

\begin{definitionbox}[The Four Paradigms]
\term{Supervised learning}: every training example is a pair $(\mathbf{x}, y)$ of
features and the correct answer (\term{label}). If $y$ is a category the task is
\term{classification}; if $y$ is a number it is \term{regression}.\\[3pt]
\term{Unsupervised learning}: examples $\mathbf{x}$ with no labels. The task is
to find structure --- groups, directions of variation, a compact
representation.\\[3pt]
\term{Semi-supervised learning}: a few labelled examples and many unlabelled
ones, used together.\\[3pt]
\term{Reinforcement learning}: an agent acts in an environment and receives a
reward. Nobody says what the right action was; the agent must discover which
actions lead, eventually, to high reward.
\end{definitionbox}

\Needspace{14\baselineskip}
\rowcolors{2}{white}{lightbg}
\begin{longtable}{@{}L{2.6cm}L{3.0cm}L{2.9cm}L{3.9cm}L{1.6cm}@{}}
\toprule
\rowcolor{primary!15}
\textbf{Paradigm} & \textbf{Experience} & \textbf{Learns} & \textbf{Example} & \textbf{Chapters} \\
\midrule
\endhead
Supervised, regression & Features and a numeric target & A function to a number &
Hours a user story will take & 5, 7, 11, 13, 14 \\
Supervised, classification & Features and a category & A function to a class &
Phishing or legitimate & 6, 8--14 \\
Unsupervised & Features only & Groups, directions, maps & Kinds of network
traffic & 15, 16 \\
Semi-supervised & Few labels, many unlabelled & A classifier using both &
Malware families from a handful of analysed samples & 17 \\
Reinforcement & States, actions, rewards & A policy: what to do in each state &
A thermostat learning to save energy & 19, 20 \\
\bottomrule
\end{longtable}

\begin{alertbox}[Classification or Regression? Look at the Target, Not the Features]
The paradigm is decided by what $y$ is. ``Predict a student's final mark'' is
regression; ``predict whether the student will pass'' is classification ---
same students, same features, different $y$. Numbers used as category codes
(port 22, grade 3 on a 1--5 scale) are a common trap: ask whether averaging two
values of $y$ would mean anything.
\end{alertbox}

\section{Designing a Learning System}

Mitchell's set text frames the construction of any learner as four decisions,
made in order. They are shown here for a problem this course returns to
repeatedly: predicting in Week~4 which students will fail a course.

\begin{worked}[The Four Design Decisions]
\textbf{1. Choose the training experience.} Records of past students with their
Week~4 data and their final result. Is it representative? Only if past cohorts
resemble this one --- a change of syllabus or of intake breaks that.

\textbf{2. Choose the target function.} What exactly should be learned?
\begin{tight}
  \item Option A: $V:\ \text{student} \to \{\text{pass}, \text{fail}\}$.
  \item Option B: $V:\ \text{student} \to [0,1]$, the probability of failing.
\end{tight}
Option B is more useful: it lets staff call the riskiest students first.

\textbf{3. Choose a representation} for the approximation $\hat{V}$ of the
target. For example, a weighted sum of features:
\[
\hat{V}(\mathbf{x}) = \sigma(w_0 + w_1\, \text{attendance} + w_2\,
\text{quiz average} + w_3\, \text{logins per week})
\]
where $\sigma$ squashes any number into $[0,1]$ (\chref{logreg}). Choosing this
form rules out, for good or ill, every pattern it cannot express.

\textbf{4. Choose a learning algorithm} that sets the weights $w_0, \dots, w_3$
from the training experience --- here, gradient descent on a loss
(\chref{math}).

\smallskip
\textbf{The first three decisions matter more than the fourth.} A poor choice
of experience or target cannot be rescued by a better algorithm.
\end{worked}

\section{Learning as Search Through a Hypothesis Space}

Decision 3 deserves a closer look. Once you fix a representation, you have fixed
the set of all functions the learner could possibly produce.

\begin{definitionbox}[Hypothesis and Hypothesis Space]
A \term{hypothesis} $h$ is one candidate function from inputs to outputs --- one
particular setting of the weights, one particular decision tree. The
\term{hypothesis space} $H$ is the set of all hypotheses the representation
allows. The unknown true function the data came from is the \term{target
function} $c$ (or $f$). \term{Learning} is a search through $H$ for a
hypothesis that agrees with $c$ --- not just on the training data, but on data
not yet seen.
\end{definitionbox}

\dsfig{%
\begin{tikzpicture}[font=\scriptsize]
\draw[rounded corners=8pt, draw=primary, line width=1pt, fill=primary!4]
     (0,0) rectangle (7.2,4.0);
\node[font=\small\bfseries, text=primary, anchor=north west] at (0.15,3.9)
     {Hypothesis space $H$};
\foreach \x/\y in {0.7/0.6,1.4/2.9,2.3/1.1,3.1/3.2,4.6/0.5,5.9/3.1,6.5/1.2,
                   1.0/1.8,2.8/0.4,6.1/2.2,5.1/3.5,0.6/3.2}{
  \fill[gray!45] (\x,\y) circle (2pt);}
\draw[greenhl!70!black, line width=1pt, fill=greenhl!12]
     (4.2,2.05) ellipse (1.5cm and 0.95cm);
\foreach \x/\y in {3.4/2.2,3.9/1.6,4.3/2.6,4.8/1.9,5.2/2.4,3.8/2.8}{
  \fill[greenhl!70!black] (\x,\y) circle (2pt);}
\node[font=\tiny, text=greenhl!45!black, align=center] at (4.2,1.25)
     {consistent with\\the training data};
\node[star, star points=5, fill=accent, inner sep=1.6pt] at (4.45,2.25) {};
\node[font=\tiny\bfseries, text=accent, anchor=west] at (4.6,2.25) {$c$?};
% right side notes
\node[anchor=north west, align=left, text width=6.0cm] at (7.7,3.9)
     {\textbf{\color{primary}Grey dots:} hypotheses the representation allows but
      the data rules out.\\[4pt]
      \textbf{\color{greenhl!55!black}Green region:} every hypothesis that fits the
      training data. Usually there are many.\\[4pt]
      \textbf{\color{accent}Star:} the true target, if it is in $H$ at all.\\[6pt]
      \textbf{Which green hypothesis to return?} The data cannot say --- they all fit
      it equally well. The learner's \emph{preference} among them is its
      \term{inductive bias}.};
\end{tikzpicture}}%
{Learning as search. The data narrows the hypothesis space to those hypotheses consistent
with it; something other than the data must choose among the survivors.}{hypothesis-space}

\begin{conceptbox}[Why a Learner Must Assume Something]
Show a learner the examples $(1 \to 2)$, $(2 \to 4)$, $(3 \to 6)$ and ask for the
output at 4. ``8'' is the natural answer, but the function that returns $2x$ for
$x \le 3$ and $0$ otherwise fits the three examples just as perfectly. Nothing in
the data prefers one to the other. You prefer $2x$ because you assume the simpler
pattern continues. That assumption --- the \term{inductive bias} --- is not a
defect to be removed: a learner with no bias at all cannot generalise beyond the
examples it has seen. \chref{concept} proves this, and every later chapter names
the bias of its algorithm.
\end{conceptbox}

\section{A Short History}

\dsfig{%
\begin{tikzpicture}[font=\scriptsize,
  ev/.style={anchor=north, align=center, text width=1.68cm, font=\tiny},
  yr/.style={font=\scriptsize\bfseries, text=primary}]
\draw[primary!50, line width=1.6pt, -{Stealth[length=3mm]}] (-0.3,0) -- (15.6,0);
\foreach \x/\y/\t/\c in {
   0.3/1958/Rosenblatt's perceptron learns from examples (Ch 13)/primary,
   2.1/1959/Samuel's checkers player; the term ``machine learning''/primary,
   3.9/1986/ID3 decision trees (Ch 9); backpropagation (Ch 13)/greenhl,
   5.7/1989/Q-learning (Ch 20)/tealhl,
   7.5/1995/Support vector machines (Ch 12)/greenhl,
   9.3/1997/Boosting analysed; Mitchell's set text/greenhl,
  11.1/2001/Random forests (Ch 14)/greenhl,
  12.9/2012/Deep networks win ImageNet/purplehl,
  14.7/2016/AlphaGo: deep RL beats a world champion/tealhl}{
  \fill[\c] (\x,0) circle (3pt);
  \node[yr] at (\x,0.35) {\y};
  \node[ev] at (\x,-0.2) {\t};}
\end{tikzpicture}}%
{Selected milestones. Most of what this course teaches was in place by 2001; what changed
after 2012 was mainly scale --- data, computation, and the depth of the networks.}{history}

The ideas in this course are older than most students expect. Frank Rosenblatt's
perceptron learned from examples in 1958; Arthur Samuel's checkers program, which
improved by playing against itself, gave the field its name in 1959. Decision
trees, backpropagation, reinforcement learning, support vector machines and
ensembles were all in place by 2001. What changed after 2012 was mostly
\emph{scale} --- data, computing power, and the depth of the neural networks
that both made practical. That is why a course built on a 1997 set text is not
out of date: the deep learning course that follows this one is built on exactly
these foundations.

\section{Where Machine Learning Sits}

In the \emph{Artificial Intelligence} course you built agents that search,
reason and plan using knowledge someone supplied. Machine learning is the part
of AI in which the knowledge is \emph{acquired from data} instead. Deep learning
is, in turn, the part of machine learning that uses neural networks with many
layers. A chess engine searching with a hand-written evaluation function is AI
but not ML; the same engine with an evaluation function learned from games is
both.

\begin{pitfall}
\begin{tight}
  \item \textbf{Choosing the algorithm first.} ``Let's use a neural network''
        before $T$, $P$ and $E$ are written down. Decide what success means, and
        what data exists, before anything else.
  \item \textbf{Judging a model on the data it learned from.} Any sufficiently
        flexible model can memorise its training data. Performance means
        performance on examples it has not seen (\chref{generalization}).
  \item \textbf{Treating a category code as a number.} Port 443 is not ``more''
        than port 80; averaging them means nothing.
  \item \textbf{Believing unsupervised means objective.} Every clustering
        reflects choices --- of features, distance and number of groups --- made
        by a person (\chref{clustering}).
  \item \textbf{Calling any automation ``machine learning''.} A threshold
        written by a person is a rule, not a model, however useful it is.
\end{tight}
\end{pitfall}

%---------------------------------------------------------------------
\begin{fourlenses}{Framing a Learning Problem}
\lensLA{$T$: rank this term's students by risk of failing. $P$: of the students
who go on to fail, the fraction in the top 20\% of the ranking. $E$: three past
cohorts with their Week~4 records. The trap is $E$: a model trained on cohorts
who received no intervention learns a world that the intervention itself will
change.}

\lensPM{$T$: predict the actual hours for a user story. $P$: mean absolute error
in hours. $E$: the team's completed stories. The trap is volume --- a single
team completes perhaps two hundred stories a year, which is small data, and the
estimates the team gave were themselves influenced by earlier predictions.}

\lensIOT{$T$: classify one second of motor vibration as healthy or failing.
$P$: failures caught at least a day before breakdown, with at most one false
alarm per machine per month. $E$: vibration logs joined to maintenance records.
The trap is that failures are rare, so there may be only a dozen labelled
failures in years of data.}

\lensSEC{$T$: classify each network connection as benign or malicious. $P$:
detection rate at a fixed, tolerable number of alerts per analyst per hour.
$E$: labelled traffic from past incidents. The trap is the adversary: $E$
describes last year's attacks, and the attacker's whole aim is for this year's
not to resemble them.}
\end{fourlenses}

%---------------------------------------------------------------------
\begin{lab}[A Rule Against a Learned Model]
Write a rule by hand, then let an algorithm learn one, and compare them on data
neither has seen. The Wisconsin breast cancer data ships with scikit-learn: 569
tumours, 30 measurements each, labelled malignant or benign.
\begin{lstlisting}[language=Python]
from sklearn.datasets import load_breast_cancer
from sklearn.model_selection import train_test_split
from sklearn.tree import DecisionTreeClassifier
from sklearn.metrics import accuracy_score

data = load_breast_cancer()
X, y = data.data, data.target            # y: 0 = malignant, 1 = benign
X_tr, X_te, y_tr, y_te = train_test_split(
    X, y, test_size=0.25, random_state=0, stratify=y)

# --- 1. a rule written by a person: large tumours are malignant ---------
radius = list(data.feature_names).index("mean radius")
rule = (X_te[:, radius] < 15).astype(int)     # small -> benign (1)
print("hand-written rule :", accuracy_score(y_te, rule))

# --- 2. a rule found by a learning algorithm ----------------------------
tree = DecisionTreeClassifier(max_depth=3, random_state=0).fit(X_tr, y_tr)
print("learned tree      :", accuracy_score(y_te, tree.predict(X_te)))

# --- 3. the baseline every result must beat -----------------------------
majority = int(y_tr.mean() > 0.5)
print("always 'majority' :", accuracy_score(y_te, [majority] * len(y_te)))
\end{lstlisting}
Then change the threshold 15 in the hand-written rule, and the depth 3 of the
tree, and watch both numbers. Which one is easier to improve by hand?
\end{lab}

%---------------------------------------------------------------------
\begin{checkpoint}
\begin{tightnum}
  \item A thermostat switches on the heating whenever the temperature drops below
        a value the owner set. Is this machine learning? Use the definition.
  \item Classify each as regression, classification, clustering or reinforcement
        learning: (a)~predicting tomorrow's server load in requests per second;
        (b)~grouping customers with no labels; (c)~deciding whether a login is
        fraudulent; (d)~a program learning to play a video game from its score.
  \item Why can a learner with no inductive bias not classify an example it has
        never seen?
\end{tightnum}
\end{checkpoint}

%---------------------------------------------------------------------
\begin{chaptersummary}
\begin{tight}
  \item Machine learning reverses traditional programming: you supply examples
        of the answers and the algorithm produces the rules.
  \item A program learns if its performance $P$ at task $T$ improves with
        experience $E$. Name all three before choosing an algorithm.
  \item The paradigm is set by the experience: labelled (supervised), unlabelled
        (unsupervised), partly labelled (semi-supervised) or rewarded
        (reinforcement). Regression and classification differ only in $y$.
  \item Designing a learner means choosing the experience, the target function,
        its representation and a learning algorithm --- in that order of
        importance.
  \item Learning is search through a hypothesis space. The data narrows it; the
        inductive bias chooses among what is left, and without one no
        generalisation is possible.
\end{tight}
\end{chaptersummary}

\begin{reviewq}
  \item \textbf{(MCQ)} Predicting the number of help-desk tickets tomorrow is:
        (a)~classification (b)~regression (c)~clustering (d)~reinforcement
        learning.
  \item \textbf{(MCQ)} In Mitchell's definition, $E$ stands for: (a)~error
        (b)~experience (c)~expectation (d)~evaluation.
  \item \textbf{(MCQ)} A dataset of 10{,}000 photographs of which 200 are labelled
        suits: (a)~supervised (b)~unsupervised (c)~semi-supervised
        (d)~reinforcement learning.
  \item \textbf{(Short)} Give one problem for which you would write rules rather
        than learn a model, and explain why.
  \item \textbf{(Short)} Explain, with an example, the difference between the
        target function and a hypothesis.
  \item \textbf{(Short)} Why is ``accuracy'' often the wrong performance measure?
        Give a numeric example.
  \item \textbf{(Applied)} A university wants to recommend elective courses to
        students. Specify $T$, $P$ and $E$; state which paradigm applies; and
        describe the four design decisions you would make.
  \item \textbf{(Applied)} For the phishing filter in the first worked example,
        describe what could make the training experience unrepresentative of the
        e-mail the filter will face next year, and what you would do about it.
\end{reviewq}
